
\section{Distribuições}

\subsection*{Funções testes}
\begin{frame}{Funções Testes}
	\begin{defin}
		Definimos o \dt{espaço das funções testes}, denotado por $\D(\Omega)$, o conjunto $C_c^\infty(\Omega)$ dotado da seguinte noção de convergência: Dados uma sequência $(\varphi_m)$ em  $\D(\Omega)$ e $\varphi\in \D(\Omega)$, dizemos que 
		\[\varphi_m\to \varphi \text{ em } \D(\Omega),\]
		quando existe $K\subset \Omega$ compacto tal que 
		\begin{enumerate}[$i$]
			\item $\supp \varphi_m$, $\supp \varphi\subset K$, $\forall m\in \N$;
			\item $D^\alpha \varphi_m\to D^\alpha \varphi$ uniformemente em $K$, $\forall \alpha\in \N^n$.
		\end{enumerate}
	\end{defin}


	Na verdade, $\D(\Omega)$ é o espaço vetorial $C_c^\infty(\Omega)$ munido de uma topologia definida por uma família de seminormas. Esta topologia nos dá a noção de convergência  acima. Para mais detalhes, veja \cite[Appendix 2]{kesavan} ou \cite[Seção 1.3]{cavalcanti-dist}.

\end{frame}



\subsection*{Distribuições}
\begin{frame}{Distribuições sobre $\Omega$}
	\begin{defin}
		Seja $\Omega$ um aberto de $\R^\dm$.	Uma \dt{distribuição sobre $\Omega$} é um funcional linear $T:\D(\Omega)\longrightarrow \R$ que é contínua no sentido da convergência sobre $\D(\Omega)$, isto é, 
		\begin{center}
			para toda $\varphi_m\to 0$ em $\D(\Omega)$, tem-se que 
			$\inp{T}{\varphi_m}\to 0$.
		\end{center}
		
		
		O espaço das distribuições é o dual do espaço das funções testes $\D(\Omega)$ e é denotado por $\D'(\Omega)$.
	\end{defin}
\end{frame}

\begin{frame}
		\begin{exe}
		 	Se $u\in \Lloc^1(\Omega)$, então $T_u$ definida por 
				\[\inp{T_u}{\varphi}=\int_\Omega u(x)\varphi(x)\,dx, \ \varphi \in \D(\Omega)\]
				é uma distribuição sobre $\Omega$.
	
	\end{exe}
	\begin{prop}
	A aplicação $T: u\in \Lloc^1(\Omega)\longmapsto T_u \in \D'(\Omega)$ é linear, injetiva e contínua, isto é, $\Lloc^1(\Omega)\hookrightarrow \D'(\Omega)$. 	
\end{prop}

\end{frame}

\subsection*{Distribuição delta de Dirac}
\begin{frame}{Distribuição delta de Dirac}
	\begin{example}
		\begin{enumerate}
			\item 		Dado $x_0\in \Omega $, a aplicação $\delta_{x_0}$ definida por 
			\[\inp{\delta_{x_0}}{\varphi}=\varphi(x_0),\ \varphi\in \D(\Omega),\]
			é uma distribuição sobre $\Omega$ denominada \dt{delta de Dirac}\footnote{$\delta_0$ é conhecida por ``função $\delta$ de Dirac" e foi introduzida por Paul Dirac em 1930. Seus cálculos simbólicos não faziam sentido matemático até a correta formulação pela Teoria das Distribuições desenvolvida por Laurent Schwartz.}.
			
			\item A distribuição $\delta_{x_0}$ não é definida por uma função localmente integrável.
		\end{enumerate}
		
	\end{example}

\end{frame}

\begin{frame}
	\begin{exer}
		Mostre que são distribuições:
\begin{enumerate}
	\item Seja $\lambda $ um caminho retificável de classe $C^1([a,b])$ cuja imagem está contida em $\Omega$. Defina
	\[\inp{\delta_\lambda}{\varphi}=\int_\lambda \varphi\, ds,\ \forall \varphi\in \D(\Omega).\]
	
	\item Seja $v(x)=1/x,$ $x\neq 0$. Defina
	\[\inp{VP(1/x)}{\varphi}=\lim\limits_{\ep\to 0}\int_{|x|>\ep}\frac{\varphi}{x}\,dx, \ \forall \varphi\in \D(\R),\]
	denominada distribuição Valor Principal de $1/x$.
	
	\item O funcional $\varphi \in D(\Omega) \longmapsto \varphi'(0)\in \R$, conhecido como doublet ou dipole distribuition.
\end{enumerate}
	\end{exer}
\subsection*{Exercício \theexer}
\end{frame}


\subsection*{Convergência de Distribuições}

\begin{frame}{Convergência de Distribuições}
	Sendo  $\D'(\Omega)$ o  dual de um espaço vetorial topológico, podemos considerar sobre ele a topologia fraca*. Com essa topologia, temos a seguinte noção de convergência.
	
	\begin{defin}
Dizemos que $T_n\to T$ em $\D'(\Omega)$ se
\[\inp{T_n}{\varphi}\to \inp{T}{\varphi},\ \forall \varphi \in \D(\Omega).\]
	\end{defin}
	\begin{exe}
A sequência de funções regularizantes $(\rho_n)$ definidas no Exemplo \ref{seq-regul} converge para a distribuição delta de Dirac $\delta_{0}$. Isto é, uma sequência de funções localmente integráveis pode convergir para um distribuição que não é definida por uma função localmente integrável.
	\end{exe}
\end{frame}

\subsection*{Derivadas das Distribuições}
\begin{frame}{Derivadas das Distribuições}
\begin{defin}
	Sejam $T\in \D'(\Omega)$ e $\alpha\in \N^\dm$. A \dt{derivada} de ordem $\alpha$ de $T$ é a distribuição $D^\alpha T$ definida por:
	\[\inp{D^\alpha T}{\varphi}=(-1)^{|\alpha|}\inp{T}{D^\alpha \varphi},\ \forall \varphi \in \D(\Omega).\]
\end{defin}

\begin{exe}
	\begin{enumerate}[a]
		\item Se $u\in C^1(\Omega)$ então a derivada no sentido das distribuições coincide com a derivada no sentido clássico.
		
		\item Se $\Omega=(-1,1)$ então a função  
		$u(x)=\chi_{(0,1)}$ é localmente integrável, com derivada clássica $u'=0$, exceto em $x=0$, mas $Du=\delta_0$ no sentido das distribuições.
		
	\end{enumerate}
\end{exe}
\end{frame}

	
\subsection*{Derivada Fraca}
\begin{frame}
	
\begin{defin}
	Dados  $u\in \Lloc^1(\Omega)$ e $\alpha\in \N^\dm$, dizemos que $D^\alpha u\in \D'(\Omega)$ é a \dt{derivada fraca} de $u$ quando $D^\alpha u\in \Lloc^1(\Omega)$, isto é, quando existe $v\in \Lloc^1(\Omega)$ tal que 
	\[\inp{D^\alpha u}{\varphi}=\int_\Omega v\varphi\, dx, \forall \varphi\in \D(\Omega).\]
\end{defin}

\begin{prop}
Se $f$ é absolutamente contínua em $[a,b]$, então $f$ é derivável q.s. em $[a,b]$ e as derivadas clássicas e distribucional coincidem.
\end{prop}

\begin{prop}
	Se $T_n\to T$ em $\D'(\Omega)$, então $D^\alpha T_n\to D^\alpha T$ em $\D'(\Omega)$, para qualquer multi-índice $\alpha$. Em outras palavras, o operador $D^\alpha: \D'(\Omega) \longrightarrow \D'(\Omega)$ é linear e contínuo.
\end{prop}
\end{frame}

\subsection*{Multiplicação de distribuições por funções $C^\infty$}
\begin{frame}{Multiplicação de distribuições por funções $C^\infty$}
\begin{defin}
	Sejam $T\in \D'(\Omega)$ e $\psi\in C^\infty(\Omega) $. O \dt{produto $\psi T$} por:
	\[\inp{\psi T}{\varphi}=\inp{T}{\psi \varphi},\ \forall \varphi\in \D(\Omega) \]
	define uma distribuição sobre $\Omega$.
\end{defin}

\begin{exer}
	Prove a fórmula de Leibniz para distribuições: 	Seja $\Omega\subset \R^\dm$ um aberto. Se $\psi\in C^\infty(\Omega)$ e $T\in \D'(\Omega)$, então
	\[D^\alpha(\psi T)=\sum_{\beta\leq \alpha}\frac{\alpha!}{\beta!(\alpha-\beta)!}D^\beta\psi D^{\alpha-\beta}T,\ \forall \alpha\in \N^\dm.\]
\end{exer}
\subsection*{Exercício \theexer}
\end{frame}

